\documentclass[a4paper,fontset=windows]{ctexart}
\usepackage{fontspec}
\usepackage{color}
\definecolor{dkgreen}{rgb}{0,0.5,0}
\definecolor{dkred}{rgb}{0.5,0,0}
\usepackage{graphicx}
\setsansfont{Consolas}
\setmonofont{Consolas}
\usepackage{geometry}
\geometry{top=3cm,bottom=3cm,left=2.5cm,right=2.5cm}

\usepackage{listings}
\lstset{
	language=Python,
	basicstyle=\ttfamily\small,
	aboveskip={1.0\baselineskip},
	belowskip={1.0\baselineskip},
	columns=fixed,
	extendedchars=true,
	breaklines=true,
	tabsize=4,
	prebreak=\raisebox{0ex}[0ex][0ex]{\ensuremath{\hookleftarrow}},
	frame=lines,
	showtabs=false,
	showspaces=false,
	showstringspaces=false,
	keywordstyle=\color[rgb]{0,0.5,0},
	commentstyle=\color[rgb]{0.5,0,0},
	stringstyle=\color[rgb]{1,0,0},
	numbers=left,
	numberstyle=\small\ttfamily ,
	stepnumber=1,
	numbersep=12pt,
	captionpos=t,
	escapeinside={\%*}{*)}
	}

\begin{document}
	
	\title{\textbf{数据结构与算法K06\ \ \ \ 课堂练习\\\vspace{1ex}单词最小编辑距离的递归解法}}
	\author{张植竣\ 1900011604\ 北京大学物理学院天文学系}
	\date{2020年3月31日}
	
	\maketitle
	
	\section*{源代码}

	\begin{lstlisting}
	def dpWordEdit(original, target, oplist):
		score = 0
		operations = []
		
		def Insert(s):
			return oplist['insert']  # 添加一个字母
		
		def Delete(t):
			return oplist['delete']  # 删除一个字母
		
		def Replace(m, n):
			if m == n:
				return oplist['copy']  # 替换一个相同的字母 == 复制一个字母
			else:
				return oplist['insert'] + oplist['delete']  # 替换一个不同的字母 == 添加一个字母，再删除一个字母
		
		originalList = list(original)  # 源单词列表
		targetList = list(target)  # 目标单词列表
		
		MinScore = {}
		# 字典，储存递归过程中的中间步骤
		solution = {(i, j):[] for i in range(len(originalList) + 1)
							  for j in range(len(targetList) + 1)}
		# 二维字典，储存由源单词中前i个字母变为目标单词中前j个字母所需要的最少操作过程
		
		def minscore(i, j):  # 双重递归，在MinScore中填写一个二维递归字典
		
			if i == 0 and j != 0:
				sco = j * oplist['insert']
				MinScore[(i, j)] = sco
				solution[(i, j)] = solution[(i, j - 1)][:]
				solution[(i, j)].append('insert {0}'.format(targetList[j - 1]))
				return sco
			# 递归字典的首行
		
			elif j == 0 and i != 0:
				sco = i * oplist['delete']
				MinScore[(i, j)] = sco
				solution[(i, j)] = solution[(i - 1, j)][:]
				solution[(i, j)].append('delete {0}'.format(originalList[i - 1]))
				return sco
			# 递归字典的首列
		
			elif i == 0 and j == 0:
				sco = 0
				MinScore[(i, j)] = sco
				solution[(i, j)] = []
				return sco
			# 递归字典的第一个元素
			
			elif (i, j) in MinScore:  # (i, j)元的值已经被储存在MinScore中，则直接返回
				sco = MinScore[(i, j)]
				return sco
		
			else:
				sco = min(minscore(i, j - 1) + Insert(targetList[j - 1]),
						  # 源单词比目标单词少一个字母，添加该字母，记录分值
						  minscore(i - 1, j - 1) + Replace(originalList[i - 1], targetList[j - 1]),
						  # 源单词与目标单词同时增加一个字母，比较这两个字母是否相同，
						  # 选择是复制该字母或是添加目标字母，再删除原字母，记录相应分值
						  minscore(i - 1, j) + Delete(originalList[i - 1]))
						  # 源单词比目标单词多一个字母，删去该字母，记录分值
		
				if sco == minscore(i, j - 1) + Insert(targetList[j - 1]):
					solution[(i, j)] = solution[(i, j - 1)][:]
					solution[(i, j)].append('insert {0}'.format(targetList[j - 1]))
					# 源单词比目标单词少一个字母，添加该字母，记录操作
		
				elif sco == minscore(i - 1, j - 1) + Replace(originalList[i - 1], targetList[j - 1]):
					solution[(i, j)] = solution[(i - 1, j - 1)][:]
					if originalList[i - 1] == targetList[j - 1]:
						solution[(i, j)].append('copy {0}'.format(targetList[j - 1]))
						# 源单词与目标单词同时增加一个字母，两个字母相同，复制该字母，记录操作
					else:
						solution[(i, j)].append('insert {0}'.format(targetList[j - 1]))
						solution[(i, j)].append('delete {0}'.format(originalList[i - 1]))
						# 源单词与目标单词同时增加一个字母，两个字母不同，添加目标字母，再删除原字母，记录操作
		
				else:
					solution[(i, j)] = solution[(i - 1, j)][:]
					solution[(i, j)].append('delete {0}'.format(originalList[i - 1]))
					# 源单词比目标单词多一个字母，删去该字母，记录操作
		
				MinScore[(i, j)] = sco
				return sco
		
		score = minscore(len(originalList), len(targetList))  # 最终分值
		operations = solution[(len(originalList), len(targetList))]  # 最终操作,即solution表中的最后一个元素
		return score, operations
		
	# 检验
	print("========= 单词最小编辑距离问题 =========")
	oplist = [{'copy': 5, 'delete': 20, 'insert': 20},
			  {'copy':5, 'delete':10, 'insert':15},
			  {'copy':10, 'delete':25, 'insert':20}]
	originalWords = ["cane", "sheep", "algorithm", "debug", "difficult", "directory", "wonderful"]
	targetWords = ["new", "sleep", "alligator", "release", "sniffing", "framework", "terrific"]
		
	import time
	a = time.time()
	for i in range(len(oplist)):  # 可以同时对三种分值设置进行检验
		for j in range(len(originalWords)):
			score, operations = dpWordEdit(originalWords[j], targetWords[j], oplist[i])
			print(score)
			print(operations)
		print('')
	b = time.time()
	print('Total time used:%.4fs'%(b-a))  # 计时程序
	\end{lstlisting}
	
	\section*{输出结果截图}

\begin{figure}[ht]
	\centering
	\includegraphics[width=1\linewidth]{1900011604_K06_picture.pdf}
\end{figure}

\end{document}